\documentclass[12pt, a4paper]{article}
\usepackage[utf8]{inputenc}
\usepackage[german]{babel}
\usepackage{lmodern}
\usepackage{amsmath, amssymb, amsthm}
\usepackage{mathtools}
\usepackage{bm}
\usepackage{graphicx}
\usepackage{xcolor}
\usepackage{tikz}
\usetikzlibrary{arrows.meta, positioning, shapes.geometric, calc,
                decorations.pathreplacing, matrix, fit, backgrounds}
\usepackage{booktabs}
\usepackage{array}
\usepackage{longtable}
\usepackage{geometry}
\usepackage{setspace}
\usepackage{parskip}
\usepackage{hyperref}
\usepackage{cleveref}
\usepackage{natbib}
\usepackage{caption}
\usepackage{subcaption}
\usepackage{algorithm}
\usepackage{algpseudocode}
\usepackage{enumitem}
\usepackage{multirow}
\usepackage{float}
\usepackage{rotating}
\usepackage{microtype}
\setlist[itemize,1]{label=--}

\emergencystretch=4em

\geometry{margin=2.5cm}

% ---------- Farben ----------
\definecolor{mainblue}{RGB}{25, 70, 140}
\definecolor{accentred}{RGB}{180, 50, 30}
\definecolor{lightgray}{RGB}{245, 245, 248}
\definecolor{darkgray}{RGB}{60, 60, 70}
\definecolor{darkgreen}{RGB}{30, 100, 50}
\definecolor{darkorange}{RGB}{200, 100, 0}

% ---------- Theorem-Umgebungen ----------
\newtheorem{theorem}{Theorem}[section]
\newtheorem{lemma}[theorem]{Lemma}
\newtheorem{proposition}[theorem]{Proposition}
\newtheorem{corollary}[theorem]{Korollar}
\theoremstyle{definition}
\newtheorem{definition}[theorem]{Definition}
\newtheorem{example}[theorem]{Beispiel}
\theoremstyle{remark}
\newtheorem{remark}[theorem]{Bemerkung}

% ---------- Hyperref ----------
\hypersetup{
  colorlinks=true,
  linkcolor=mainblue,
  citecolor=accentred,
  urlcolor=mainblue,
  pdftitle={Von der Logik: Ordnung, Sinn und Wahrheit},
  pdfauthor={Stephan Epp},
}

% ---------- Makros ----------
\newcommand{\W}{\mathrm{W}}
\newcommand{\F}{\mathrm{F}}
\newcommand{\N}{\mathbb{N}}
\newcommand{\WFF}{\mathrm{WFF}}
\newcommand{\Var}{\mathrm{Var}}
\newcommand{\Log}{\mathrm{Log}}
\newcommand{\Bel}{\mathrm{Bel}}
\newcommand{\Taut}{\mathrm{Taut}}
\newcommand{\Kon}{\mathrm{Kon}}
\newcommand{\Ord}{\mathrm{Ord}}
\newcommand{\Sinn}{\mathrm{Sinn}}
\newcommand{\len}{\ell}

\title{\textbf{Von der Logik}\\[0.5em]
  {\Large\itshape Ordnung, Sinn und Wahrheit\\Eine formale Grundlegung der Aussagenlogik}}
\author{Stephan Epp}
\date{09. September 2026}

% ============================================================
\begin{document}
% ============================================================

\maketitle
\tableofcontents
\newpage

% ============================================================
\section{Einleitung}
\label{sec:einleitung}
% ============================================================

\subsection{Die Ausgangsfrage}
\label{subsec:ausgangsfrage}

Woher kommt sie? Wer gebraucht sie? Und wohin führt sie? Ist Logik mehr als geordnete Worte, die Sinn machen? Angenommen, sie ist mehr, was muss hinzugefügt werden? Was könnte ergänzt werden, damit Logik erst wirklich Logik ist? Sind es Formeln? Sind es Zahlen? Oder sind es viele Soldaten? Nein. Denn jeder weiß, beim Erfassen des Sinns für Formeln, Zahlen oder Soldaten erfolgt dies auch in geordneten Worten mit Sinn. Also das ist es nicht. Worte allein reichen nicht. Es muss an der Ordnung und an dem Sinn liegen. Was ist Ordnung? Was ist Sinn? Ordnung ist klar. Aber was ist Sinn? Oder anders gefragt, was ist sinnvoll? Sinn hat es, wenn es sinnvoll ist. Voller Sinn oder anders, voller Wahrheit ist es dann, wenn alles nur wahr ist. Was ist wahr? Alles das, was nicht falsch ist. Vorausgesetzt, es gibt nur wahr und falsch. Logik ist also nicht falsch. Logik ist deshalb nur wahr in geordneten Worten. Logisch?

Diese Fragen sind keine rhetorischen Verzierungen, sondern der eigentliche Forschungsgegenstand dieser Arbeit. Sie enthalten -- in vorformaler Sprache -- bereits die vollständige Architektur der klassischen Aussagenlogik: einen Begriff von \emph{Ordnung} (geordnete Worte, d.\,h. eine Grammatik), einen Begriff von \emph{Sinn} (Sinnhaftigkeit, d.\,h. eine Semantik), einen Begriff von \emph{Wahrheit} als Gegenstück zu \emph{Falschheit} unter Ausschluss eines Dritten (das Bivalenzprinzip), und schließlich die Vermutung, dass Logik genau dort beginnt, wo Ordnung und Sinn zusammenfallen. Die vorliegende Arbeit nimmt diese Vermutung wörtlich und weist sie formal-mathematisch nach.

\subsection{Wissenschaftliche Einordnung}
\label{subsec:einordnung}

Die vorliegende Arbeit behandelt drei zentrale Themenkomplexe:

\begin{enumerate}[label=(\roman*)]
  \item \textbf{Ordnung als Syntax:} Formaler Nachweis, dass sich der intuitive Begriff der ``geordneten Worte'' präzise als Grammatik einer formalen Sprache (die Menge der wohlgeformten Formeln $\WFF$) fassen lässt, samt Beweis der eindeutigen Lesbarkeit jeder Formel.

  \item \textbf{Sinn als Semantik:} Rigoroser Beweis, dass der intuitive Begriff des ``Sinns'' bzw. der ``Sinnhaftigkeit'' präzise als Belegungsfunktion mit Werten in einer zweielementigen Wahrheitswertmenge $\{\W, \F\}$ gefasst werden kann (Bivalenzprinzip), und dass jede wohlgeformte Formel unter jeder Belegung genau einen Wahrheitswert besitzt (Bivalenzsatz).

  \item \textbf{Logik als Deckung von Ordnung und Sinn:} Formaler Beweis, dass ``Logik'' im engeren Sinn -- die Menge der allgemeingültigen (tautologischen) Aussagen -- eine echte Teilmenge der bloß wohlgeformten Ausdrücke ist, dass also Ordnung allein nicht hinreicht, und dass umgekehrt jede syntaktisch geordnete, in endlich vielen Schritten deduktiv abgeleitete Formel notwendig auch semantisch wahr ist und umgekehrt (Vollständigkeits- und Korrektheitssatz). Dies ist die formale Kaschierung des Schlusssatzes ``Logik ist deshalb nur wahr in geordneten Worten''.
\end{enumerate}

\subsection{Aufbau der Arbeit}
\label{subsec:aufbau}

Abschnitt~\ref{sec:syntax} legt die formale Sprache der Aussagenlogik fest und beweist die Eindeutigkeit ihrer Ordnung. Abschnitt~\ref{sec:semantik} entwickelt die Semantik (den Sinn) über Belegungen und beweist den Bivalenzsatz. Abschnitt~\ref{sec:zusammenspiel} zeigt formal, dass weder Ordnung allein noch Sinn allein hinreichen, und beantwortet die Frage, warum es weder auf Formeln, Zahlen noch ``Soldaten'' als Träger ankommt. Abschnitt~\ref{sec:kalkuel} führt einen Hilbert-Kalkül als Modell der Deduktion (der geordneten Kette von Schritten) ein und beweist den Korrektheitssatz. Abschnitt~\ref{sec:vollstaendigkeit} beweist den Vollständigkeitssatz und leitet daraus den Hauptsatz dieser Arbeit her, der Ordnung, Sinn und Logik formal in eins setzt. Abschnitt~\ref{sec:rueckbezug} übersetzt jede Frage der Einleitung explizit in ihre formale Antwort. Abschnitt~\ref{sec:ausblick} schließt mit Diskussion und Ausblick.

% ============================================================
\section{Ordnung: Die formale Sprache der Aussagenlogik}
\label{sec:syntax}
% ============================================================

\subsection{Alphabet und Grammatik}
\label{subsec:alphabet}

\begin{definition}[Alphabet]
\label{def:alphabet}
Ein \emph{aussagenlogisches Alphabet} ist eine Menge
\[
  \Sigma \;=\; \Var \;\cup\; \{\neg, \to\} \;\cup\; \{(, )\},
\]
wobei $\Var = \{p_1, p_2, p_3, \dots\}$ eine abzählbare Menge von \emph{Aussagevariablen} ist, disjunkt von den \emph{Junktoren} $\neg$ (Negation) und $\to$ (Implikation) sowie den Hilfszeichen $($ und $)$.
\end{definition}

\begin{remark}
\label{rem:soldaten}
Die konkrete Wahl der Symbole in $\Var$ ist gleichgültig. Man könnte statt $p_1, p_2, \dots$ ebenso gut Zahlen, Buchstaben eines beliebigen Alphabets, Farben oder -- um die Eingangsfrage aufzugreifen -- eine durchnummerierte Kompanie von Soldaten $s_1, s_2, \dots$ als Variablennamen verwenden. Dass dies keinen Unterschied macht, wird in Theorem~\ref{thm:isomorphie} bewiesen.
\end{remark}

\begin{definition}[Wörter über $\Sigma$]
\label{def:woerter}
Ein \emph{Wort} über $\Sigma$ ist eine endliche Folge $w = a_1 a_2 \cdots a_n$ mit $a_i \in \Sigma$ für $1 \le i \le n$ und $n \in \N_0$. Die Menge aller Wörter über $\Sigma$ wird mit $\Sigma^\ast$ bezeichnet, ihre Länge mit $\len(w) = n$.
\end{definition}

Nicht jedes Wort über $\Sigma$ ist bereits sinnvoll geordnet -- ``$)\neg\neg(p_1\to{}$'' ist ein Wort, aber Unsinn. Die \emph{Ordnung}, nach der die Einleitung fragt, ist genau die Auszeichnung derjenigen Wörter, die einer festen Bildungsregel folgen.

\begin{definition}[Wohlgeformte Formel, $\WFF$]
\label{def:wff}
Die Menge $\WFF \subseteq \Sigma^\ast$ der \emph{wohlgeformten Formeln} ist die kleinste Teilmenge von $\Sigma^\ast$ mit:
\begin{enumerate}[label=(\alph*)]
  \item $p \in \WFF$ für jedes $p \in \Var$ \hfill (Basis),
  \item wenn $\varphi \in \WFF$, dann auch $\neg\varphi \in \WFF$ \hfill (Negationsregel),
  \item wenn $\varphi, \psi \in \WFF$, dann auch $(\varphi \to \psi) \in \WFF$ \hfill (Implikationsregel).
\end{enumerate}
Formeln, die nicht durch endlich viele Anwendungen von (a)--(c) erzeugt werden können, gehören nicht zu $\WFF$. Wir schreiben $\varphi, \psi, \chi, \dots$ für beliebige Elemente von $\WFF$.
\end{definition}

\begin{remark}[Ordnung als Minimalitätsforderung]
\label{rem:minimalitaet}
Die Formulierung ``kleinste Teilmenge, die (a)--(c) erfüllt'' ist entscheidend: Sie schließt aus, dass beliebige zusätzliche Wörter willkürlich zu $\WFF$ hinzugefügt werden. Formal ist $\WFF$ der Schnitt aller Teilmengen $S \subseteq \Sigma^\ast$, die (a)--(c) erfüllen. Dies ist die exakte Bedeutung von ``Ordnung ist klar'' aus der Einleitung: Ordnung ist nichts anderes als Zugehörigkeit zu dieser eindeutig bestimmten, induktiv definierten Menge.
\end{remark}

\begin{definition}[Abkürzungen]
\label{def:abkuerzungen}
Für $\varphi, \psi \in \WFF$ definieren wir die folgenden Abkürzungen, die selbst nicht Teil des primitiven Alphabets sind, sondern Metasprache:
\[
  (\varphi \lor \psi) :\equiv (\neg\varphi \to \psi), \qquad
  (\varphi \land \psi) :\equiv \neg(\varphi \to \neg\psi), \qquad
  (\varphi \leftrightarrow \psi) :\equiv \big((\varphi \to \psi) \land (\psi \to \varphi)\big).
\]
\end{definition}

\subsection{Eindeutige Lesbarkeit}
\label{subsec:eindeutige_lesbarkeit}

Damit ``Ordnung'' als Begriff trägt, muss jede wohlgeformte Formel auf genau eine Weise nach Definition~\ref{def:wff} entstanden sein -- andernfalls wäre ihr syntaktischer Aufbau mehrdeutig und jede darauf aufbauende Semantik ill-definiert.

\begin{theorem}[Eindeutige Lesbarkeit]
\label{thm:eindeutige_lesbarkeit}
Jede Formel $\varphi \in \WFF$ ist auf genau eine der folgenden drei Weisen aufgebaut:
\begin{enumerate}[label=(\arabic*)]
  \item $\varphi = p$ für genau ein $p \in \Var$;
  \item $\varphi = \neg\varphi'$ für genau ein $\varphi' \in \WFF$;
  \item $\varphi = (\varphi_1 \to \varphi_2)$ für genau ein Paar $\varphi_1, \varphi_2 \in \WFF$.
\end{enumerate}
Insbesondere sind diese drei Fälle paarweise disjunkt (kein $\varphi \in \WFF$ fällt unter zwei der Fälle), und im Fall (3) sind $\varphi_1, \varphi_2$ eindeutig durch $\varphi$ bestimmt.
\end{theorem}

\begin{proof}
\emph{Existenz} der Zerlegung (irgendeine der drei Formen liegt vor) folgt unmittelbar aus Definition~\ref{def:wff}, da $\WFF$ durch genau die Regeln (a)--(c) erzeugt wird.

\emph{Disjunktheit der Fälle.} Wir zeigen: Jedes $\varphi \in \WFF$ beginnt entweder mit einem Symbol aus $\Var$, mit $\neg$, oder mit $($, und diese drei Anfangssymbole schließen sich aus. Wir führen dies durch strukturelle Induktion über den Aufbau von $\varphi$:
\begin{itemize}
  \item Ist $\varphi = p \in \Var$, so ist das einzige Symbol von $\varphi$ selbst aus $\Var$; kein $\Var$-Symbol ist gleich $\neg$ oder $($, da $\Sigma$ als disjunkte Vereinigung definiert wurde (Definition~\ref{def:alphabet}).
  \item Ist $\varphi = \neg\varphi'$, so beginnt $\varphi$ mit dem Symbol $\neg$.
  \item Ist $\varphi = (\varphi_1 \to \varphi_2)$, so beginnt $\varphi$ mit dem Symbol $($.
\end{itemize}
Da $\Var$, $\{\neg\}$ und $\{(\}$ paarweise disjunkte Teilmengen von $\Sigma$ sind, ist das Anfangssymbol von $\varphi$ hinreichend, um zwischen den drei Fällen zu unterscheiden; die Fälle können sich also nicht überschneiden.

\emph{Eindeutigkeit im Fall (3).} Angenommen, $(\varphi_1 \to \varphi_2) = (\psi_1 \to \psi_2)$ als Wörter in $\Sigma^\ast$, mit $\varphi_1, \varphi_2, \psi_1, \psi_2 \in \WFF$. Wir zeigen $\varphi_1 = \psi_1$ (und damit auch $\varphi_2 = \psi_2$, da die Restwörter dann übereinstimmen müssen) mittels folgendem Klammer-Zählargument: Definiere für ein Wort $w = a_1 \cdots a_n$ über $\Sigma$ die \emph{Klammerbilanz}
\[
  b(w) \;=\; \#\{i : a_i = (\} \;-\; \#\{i : a_i = )\}.
\]
Man zeigt durch strukturelle Induktion über $\WFF$, dass für jede Formel $\chi \in \WFF$ und jedes echte, nichtleere Anfangsstück $u$ von $\chi$ (d.\,h. $\chi = uv$ mit $u \ne \varepsilon$, $v \ne \varepsilon$) gilt: $b(u) \geq 0$, und $b(u) = 0$ nur dann, wenn $u$ selbst bereits eine vollständige wohlgeformte Formel ist, gefolgt von weiteren Symbolen. (Basisfall $\chi = p$: $\chi$ hat keine echten nichtleeren Anfangsstücke außer sich selbst und $b(\chi)=0$ trivial. Fall $\chi=\neg\chi'$: jedes Anfangsstück von $\chi$ ist entweder $\neg$ gefolgt von einem Anfangsstück von $\chi'$, worauf die Induktionsvoraussetzung anwendbar ist, da $\neg$ die Klammerbilanz nicht ändert. Fall $\chi=(\chi_1\to\chi_2)$: die öffnende Klammer erhöht $b$ um $1$; erst wenn das Anfangsstück ganz $\chi_1$ und den Pfeil und ganz $\chi_2$ und die schließende Klammer überstrichen hat, sinkt $b$ wieder auf $0$ -- vorher bleibt $b \ge 1$, was man wiederum durch Anwendung der Induktionsvoraussetzung auf $\chi_1$ bzw. $\chi_2$ erhält.)

Angenommen nun $\varphi_1 \ne \psi_1$ als Wörter. Da beide als Anfangsstücke desselben Wortes $\varphi_1 \varphi_2) = \psi_1 \psi_2)$ (nach Streichen der einleitenden Klammer und des mittleren Pfeils) auftreten, ist eines von beiden ein echtes Anfangsstück des anderen, etwa $\varphi_1$ echtes Anfangsstück von $\psi_1$. Da $\psi_1 \in \WFF$, folgt aus dem eben gezeigten Klammer-Lemma, angewandt auf $\chi = \psi_1$ und $u = \varphi_1$: entweder $b(\varphi_1) > 0$, oder $\varphi_1$ ist selbst schon eine vollständige Formel gefolgt von einem nichttrivialen Rest innerhalb von $\psi_1$. Da aber auch $\varphi_1 \in \WFF$ vollständig ist, muss nach demselben Lemma, angewandt auf $\chi=\varphi_1$ selbst, $b(\varphi_1) = 0$ gelten (jede vollständige Formel hat Klammerbilanz $0$, was man direkt durch strukturelle Induktion sieht). Beide Optionen widersprechen sich: Ist $b(\varphi_1) = 0$ und $\varphi_1$ echtes Anfangsstück von $\psi_1 \in \WFF$, so besagt das Lemma gerade, dass dies nur möglich ist, wenn $\varphi_1$ bereits eine vollständige Teilformel ist, die von $\psi_1$ dann durch weitere Symbole fortgesetzt wird -- was der Konstruktion von $\psi_1$ als atomarer, negierter oder einzelner Implikationsformel widerspricht, da $\psi_1 \in \WFF$ nach Induktionsannahme (Fallunterscheidung nach Theorem~\ref{thm:eindeutige_lesbarkeit}, Teil 1) keine echte wohlgeformte Teilformel als striktes Anfangsstück besitzt, ohne mit ihr identisch zu sein. Widerspruch. Also $\varphi_1 = \psi_1$, und da die Gesamtwörter übereinstimmen, folgt auch $\varphi_2 = \psi_2$.

Die Fälle (1) und (2) sind analog (und einfacher, da keine Klammerzählung nötig ist): In Fall (1) ist $\varphi$ ein Wort der Länge $1$, also ist $p$ durch $\varphi$ selbst eindeutig bestimmt. In Fall (2) ist $\varphi'$ das um das führende Symbol $\neg$ gekürzte Wort, also durch $\varphi$ eindeutig bestimmt.
\end{proof}

\begin{corollary}[Rekursionsprinzip auf $\WFF$]
\label{cor:rekursion}
Theorem~\ref{thm:eindeutige_lesbarkeit} rechtfertigt die \emph{strukturelle Rekursion}: Für jede Menge $M$, jede Funktion $f_0 : \Var \to M$ und jede Funktionen $f_\neg : M \to M$, $f_\to : M \times M \to M$ existiert genau eine Funktion $F : \WFF \to M$ mit
\[
  F(p) = f_0(p), \qquad F(\neg\varphi) = f_\neg(F(\varphi)), \qquad F((\varphi\to\psi)) = f_\to(F(\varphi), F(\psi)).
\]
\end{corollary}
\begin{proof}
Existenz folgt durch strukturelle Induktion über den Aufbau von $\WFF$ (Definition der Werte entlang der Erzeugungsregeln), Eindeutigkeit unmittelbar aus der eindeutigen Lesbarkeit (Theorem~\ref{thm:eindeutige_lesbarkeit}): Da jede Formel auf genau eine Art erzeugt ist, gibt es für $F(\varphi)$ keine widersprüchliche zweite Vorschrift.
\end{proof}

Korollar~\ref{cor:rekursion} ist das technische Rückgrat der gesamten Arbeit: Er erlaubt es, im nächsten Abschnitt den ``Sinn'' einer Formel wohldefiniert -- als rekursiv über die Ordnung von $\WFF$ definierte Funktion -- einzuführen.

% ============================================================
\section{Sinn: Belegung, Wahrheit und Bivalenz}
\label{sec:semantik}
% ============================================================

\subsection{Das Bivalenzprinzip}
\label{subsec:bivalenz}

Die Einleitung setzt voraus: ``es gibt nur wahr und falsch''. Dies ist keine ableitbare Tatsache, sondern eine \emph{Modellannahme} der klassischen Logik, die wir hier explizit als Axiom festhalten.

\begin{definition}[Wahrheitswertmenge, Bivalenzprinzip]
\label{def:bivalenzprinzip}
Sei $\mathbb{B} = \{\W, \F\}$ eine zweielementige Menge, deren Elemente \emph{wahr} ($\W$) und \emph{falsch} ($\F$) genannt werden, mit $\W \ne \F$. Das \emph{Bivalenzprinzip} (\emph{principium bivalentiae}, auch \emph{tertium non datur}) postuliert, dass jeder Aussage genau einer dieser zwei Werte zukommt -- es gibt kein Drittes. Formal: Für jede Aussagevariable existiert eine Zuordnung nach $\mathbb{B}$, und diese Zuordnung ist total und einwertig.
\end{definition}

\begin{remark}
\label{rem:bivalenz_ist_axiom}
Das Bivalenzprinzip ist \emph{nicht} beweisbar; es ist eine Setzung. Andere Logiken (dreiwertige Logiken, intuitionistische Logik, Fuzzy-Logik) verzichten auf genau dieses Axiom. Die Aussage der Einleitung ``vorausgesetzt, es gibt nur wahr und falsch'' benennt diese Setzung exakt und korrekt als Voraussetzung, nicht als Herleitung -- eine bemerkenswert präzise Selbstverortung. Die gesamte übrige Arbeit operiert unter dieser Voraussetzung (klassische, zweiwertige Aussagenlogik).
\end{remark}

\begin{definition}[Falschheit als Nicht-Wahrheit]
\label{def:falsch_ist_nicht_wahr}
Unter dem Bivalenzprinzip gilt für $x \in \mathbb{B}$:
\[
  x = \F \quad \Longleftrightarrow \quad x \ne \W.
\]
Dies formalisiert den Satz ``Was ist wahr? Alles das, was nicht falsch ist'' -- äquivalent, und ebenso beweisbar aus Definition~\ref{def:bivalenzprinzip}: $x = \W \Leftrightarrow x \ne \F$.
\end{definition}
\begin{proof}
``$\Rightarrow$'': Ist $x = \F$, so ist $x \ne \W$ wegen $\W \ne \F$ (Definition~\ref{def:bivalenzprinzip}). ``$\Leftarrow$'': Ist $x \ne \W$, so ist $x \in \mathbb{B}\setminus\{\W\} = \{\F\}$, da $\mathbb{B}=\{\W,\F\}$ per Definition nur zwei Elemente hat; also $x=\F$.
\end{proof}

\subsection{Belegungen und die Bedeutungsfunktion}
\label{subsec:belegungen}

\begin{definition}[Belegung]
\label{def:belegung}
Eine \emph{Belegung} ist eine Funktion $v : \Var \to \mathbb{B}$. Die Menge aller Belegungen wird mit $\Bel$ bezeichnet.
\end{definition}

\begin{definition}[Wahrheitswertfunktion, Sinn einer Formel unter einer Belegung]
\label{def:wahrheitswertfunktion}
Zu jeder Belegung $v \in \Bel$ definieren wir mittels der strukturellen Rekursion aus Korollar~\ref{cor:rekursion} (mit $M = \mathbb{B}$, $f_0 = v$) eine eindeutig bestimmte Fortsetzung $\hat v : \WFF \to \mathbb{B}$ durch:
\begin{align}
  \hat v(p) &= v(p), \qquad p \in \Var, \label{eq:sem_atom}\\
  \hat v(\neg\varphi) &= \begin{cases} \W, & \hat v(\varphi) = \F \\ \F, & \hat v(\varphi) = \W \end{cases}, \label{eq:sem_neg}\\
  \hat v((\varphi \to \psi)) &= \begin{cases} \F, & \hat v(\varphi) = \W \text{ und } \hat v(\psi) = \F \\ \W, & \text{sonst.} \end{cases} \label{eq:sem_impl}
\end{align}
Wir nennen $\hat v(\varphi)$ den \emph{Wahrheitswert von $\varphi$ unter $v$} und sagen, $\varphi$ \emph{habe unter $v$ Sinn im starken Sinne von ``wahr''}, wenn $\hat v(\varphi) = \W$; andernfalls, dass $\varphi$ unter $v$ falsch sei.
\end{definition}

Diese Definition beantwortet die scheinbar zirkuläre Formulierung der Einleitung, ``Sinn hat es, wenn es sinnvoll ist'': Isoliert betrachtet ist der Satz eine Tautologie der Umgangssprache und daher inhaltsleer. Formalisiert man ihn jedoch als Frage nach der \emph{Wohldefiniertheit von $\hat v$} -- gibt es zu jeder Formel unter jeder Belegung überhaupt und eindeutig einen Sinn (Wahrheitswert)? --, so wird daraus ein echtes, beweisbedürftiges und beweisbares mathematisches Theorem.

\begin{theorem}[Bivalenzsatz]
\label{thm:bivalenzsatz}
Für jede Formel $\varphi \in \WFF$ und jede Belegung $v \in \Bel$ existiert genau ein Wahrheitswert $\hat v(\varphi) \in \{\W, \F\}$.
\end{theorem}
\begin{proof}
\emph{Existenz und Eindeutigkeit} folgen unmittelbar aus Korollar~\ref{cor:rekursion}: $\hat v$ ist als die eindeutig bestimmte strukturell-rekursive Fortsetzung von $v$ definiert, und Korollar~\ref{cor:rekursion} garantiert, dass eine solche Fortsetzung existiert und eindeutig ist, \emph{vorausgesetzt}, die Rekursionsgleichungen \eqref{eq:sem_atom}--\eqref{eq:sem_impl} definieren für jeden Fall (Basis, Negation, Implikation) tatsächlich einen wohldefinierten Wert in $\mathbb{B}$. Dies ist der Fall:
\begin{itemize}
  \item In \eqref{eq:sem_atom} ist $v(p) \in \mathbb{B}$ nach Definition~\ref{def:belegung} für jedes $p$ definiert und eindeutig, da $v$ eine Funktion ist.
  \item In \eqref{eq:sem_neg} ist für jeden der zwei möglichen Werte $\hat v(\varphi) \in \{\W,\F\}$ (Induktionsvoraussetzung: $\varphi$ hat bereits genau einen Wert) genau ein Fall der Fallunterscheidung einschlägig, da $\W \ne \F$; der Ergebniswert ist in beiden Fällen eindeutig als Element von $\mathbb{B}$ angegeben.
  \item In \eqref{eq:sem_impl} deckt die Fallunterscheidung ``($\hat v(\varphi)=\W$ und $\hat v(\psi)=\F$)'' gegen ``sonst'' alle $2\times 2 = 4$ Kombinationen von $(\hat v(\varphi), \hat v(\psi)) \in \mathbb{B}\times\mathbb{B}$ widerspruchsfrei und vollständig ab, da ``sonst'' logisch-mengentheoretisch exakt das Komplement des ersten Falls in $\mathbb{B}\times\mathbb{B}$ bezeichnet.
\end{itemize}
Da somit jede der drei Rekursionsklauseln für jeden Input genau einen Output in $\mathbb{B}$ liefert, liefert Korollar~\ref{cor:rekursion} eine eindeutig bestimmte totale Funktion $\hat v: \WFF \to \mathbb{B}$. Insbesondere hat jede $\varphi \in \WFF$ unter jeder Belegung $v$ genau einen der zwei Werte $\W$ oder $\F$ -- niemals keinen, niemals mehrere.
\end{proof}

\begin{remark}
\label{rem:bivalenzsatz_bedeutung}
Theorem~\ref{thm:bivalenzsatz} ist die exakte, nicht-zirkuläre Fassung von ``Sinn hat es, wenn es sinnvoll ist'': Jede geordnete Formel (jedes $\varphi \in \WFF$) \emph{hat} unter jeder Interpretation \emph{Sinn} in dem präzisen Sinne, dass ihr eindeutig ein Wahrheitswert zukommt. Die Ordnung (Zugehörigkeit zu $\WFF$) ist dabei die notwendige Bedingung, damit die Sinnzuweisung überhaupt wohldefiniert ist -- ein Wort außerhalb von $\WFF$, etwa $)\neg\neg(p_1\to{}$, besitzt keine nach \eqref{eq:sem_atom}--\eqref{eq:sem_impl} rekursiv bestimmbare Bedeutung, weil das Rekursionsprinzip aus Korollar~\ref{cor:rekursion} nur auf $\WFF$ selbst, nicht auf beliebige Wörter aus $\Sigma^\ast$, anwendbar ist.
\end{remark}

\subsection{Allgemeingültigkeit: voller Sinn}
\label{subsec:tautologie}

\begin{definition}[Tautologie, Kontradiktion, Erfüllbarkeit]
\label{def:tautologie}
Sei $\varphi \in \WFF$.
\begin{enumerate}[label=(\alph*)]
  \item $\varphi$ heißt \emph{Tautologie} (allgemeingültig, ``von vollem Sinn''), geschrieben $\models \varphi$, wenn $\hat v(\varphi) = \W$ für \emph{jede} Belegung $v \in \Bel$.
  \item $\varphi$ heißt \emph{Kontradiktion}, wenn $\hat v(\varphi) = \F$ für jede Belegung $v \in \Bel$.
  \item $\varphi$ heißt \emph{erfüllbar}, wenn es (mindestens) eine Belegung $v$ gibt mit $\hat v(\varphi) = \W$.
\end{enumerate}
\end{definition}

Dies ist die exakte Fassung von ``Voller Sinn oder anders, voller Wahrheit ist es dann, wenn alles nur wahr ist'': ``alles'' quantifiziert über \emph{alle} Belegungen $v \in \Bel$, und ``nur wahr'' fordert $\hat v(\varphi) = \W$ ausnahmslos.

\begin{example}
\label{ex:tautologien}
Die Formel $(p \to p)$ ist eine Tautologie: Für $v(p)=\W$ ist $\hat v(p\to p) = \W$ nach \eqref{eq:sem_impl}, da nicht ($\hat v(p)=\W$ und $\hat v(p)=\F$) gilt; für $v(p)=\F$ ebenso, da $\hat v(p)=\W$ falsch ist, also der ``sonst''-Fall greift. Die Formel $(p \land \neg p)$ (in der Abkürzung von Definition~\ref{def:abkuerzungen}, also eigentlich $\neg(p\to\neg\neg p)$) ist eine Kontradiktion; die Formel $p$ selbst ist erfüllbar, aber weder Tautologie noch Kontradiktion.
\end{example}

% ============================================================
\section{Das Zusammenspiel von Ordnung und Sinn}
\label{sec:zusammenspiel}
% ============================================================

Wir zeigen nun formal, warum weder Ordnung allein noch Sinn allein für ``Logik'' hinreichen -- und warum es dabei, wie in der Einleitung vermutet, weder auf Formeln, Zahlen noch andere konkrete Träger der Ordnung ankommt.

\subsection{Ordnung allein reicht nicht}
\label{subsec:ordnung_reicht_nicht}

\begin{definition}[Logik als Menge]
\label{def:logik_menge}
Wir definieren formal
\[
  \Log \;:=\; \{\varphi \in \WFF : \; \models \varphi\} \;\subseteq\; \WFF.
\]
$\Log$ ist die Formalisierung des Wortes ``Logik'' im Schlusssatz der Einleitung: die Menge der Aussagen, die \emph{sowohl} geordnet (Element von $\WFF$) \emph{als auch} von vollem Sinn (allgemeingültig wahr) sind.
\end{definition}

\begin{theorem}[Ordnung ist nicht hinreichend für Logik]
\label{thm:ordnung_nicht_hinreichend}
$\Log \subsetneq \WFF$, das heißt, $\Log$ ist eine \emph{echte} Teilmenge von $\WFF$: Es gibt wohlgeformte Formeln, die nicht zu $\Log$ gehören.
\end{theorem}
\begin{proof}
$\Log \subseteq \WFF$ gilt nach Definition~\ref{def:logik_menge} unmittelbar. Zu zeigen bleibt $\Log \ne \WFF$, d.\,h. die Existenz eines Gegenbeispiels. Betrachte $\varphi := p_1 \in \Var \subseteq \WFF$ (nach Definition~\ref{def:wff}(a) wohlgeformt). Wähle die Belegung $v_0$ mit $v_0(p_1) = \F$ (eine solche Belegung existiert, da Definition~\ref{def:belegung} keinerlei Einschränkung an $v$ auf einzelnen Variablen macht: man setze etwa $v_0(p_1)=\F$ und $v_0(p_i)=\W$ für alle $i\ge 2$, was eine wohldefinierte Funktion $\Var\to\mathbb{B}$ ist). Dann ist $\hat v_0(p_1) = v_0(p_1) = \F \ne \W$ nach \eqref{eq:sem_atom}. Also existiert eine Belegung, unter der $p_1$ nicht wahr ist, somit $\not\models p_1$ nach Definition~\ref{def:tautologie}(a) (Negation von ``für jede Belegung wahr''), also $p_1 \in \WFF \setminus \Log$. Damit ist $\Log \ne \WFF$, und zusammen mit $\Log \subseteq \WFF$ folgt $\Log \subsetneq \WFF$.
\end{proof}

\begin{corollary}
\label{cor:ordnung_reicht_nicht}
Ordnung (Zugehörigkeit zu $\WFF$) ist eine notwendige, aber nicht hinreichende Bedingung für Logik (Zugehörigkeit zu $\Log$).
\end{corollary}
\begin{proof}
Notwendigkeit ist die Inklusion $\Log \subseteq \WFF$ aus Definition~\ref{def:logik_menge}; Nicht-Hinreichendheit ist genau die Aussage $\Log \ne \WFF$ aus Theorem~\ref{thm:ordnung_nicht_hinreichend}, angewendet auf das dort konstruierte Gegenbeispiel $p_1$, das geordnet, aber nicht logisch (nicht allgemeingültig) ist.
\end{proof}

Korollar~\ref{cor:ordnung_reicht_nicht} ist die formale Bestätigung von ``Worte allein reichen nicht'': Selbst wenn ein Ausdruck vollständig der Grammatik $\WFF$ genügt (``geordnete Worte''), muss er nicht wahr unter jeder Belegung sein.

\subsection{Sinn ohne Ordnung ist nicht definierbar}
\label{subsec:sinn_ohne_ordnung}

\begin{proposition}[Sinn setzt Ordnung voraus]
\label{prop:sinn_setzt_ordnung_voraus}
Die Bedeutungsfunktion $\hat v$ aus Definition~\ref{def:wahrheitswertfunktion} ist ausschließlich auf $\WFF$ definiert; für $w \in \Sigma^\ast \setminus \WFF$ existiert keine mit \eqref{eq:sem_atom}--\eqref{eq:sem_impl} verträgliche Fortsetzung.
\end{proposition}
\begin{proof}
Nach Korollar~\ref{cor:rekursion} beruht die Existenz und Eindeutigkeit von $\hat v$ auf der eindeutigen Lesbarkeit von $\WFF$ (Theorem~\ref{thm:eindeutige_lesbarkeit}). Für $w \notin \WFF$ ist $w$ nicht durch Regel (a), (b) oder (c) aus Definition~\ref{def:wff} erzeugt; die Rekursionsklauseln \eqref{eq:sem_atom}--\eqref{eq:sem_impl} setzen aber jeweils voraus, dass das Argument entweder atomar ist, mit $\neg$ beginnt und ein wohlgeformtes $\varphi$ enthält, oder die Form $(\varphi\to\psi)$ mit wohlgeformten $\varphi,\psi$ hat (Theorem~\ref{thm:eindeutige_lesbarkeit}). Für $w \notin \WFF$ ist keiner dieser drei Fälle erfüllt (sonst wäre $w \in \WFF$ nach Definition~\ref{def:wff} selbst), sodass keine der drei Klauseln anwendbar ist und $\hat{}(w)$ unbestimmt bleibt.
\end{proof}

Dies ist das formale Gegenstück zu Corollary~\ref{cor:ordnung_reicht_nicht}: So wie Ordnung ohne Sinn ``Logik'' verfehlt, ist Sinn (im Sinne einer wohldefinierten Bedeutungszuweisung) ohne Ordnung nicht einmal \emph{definiert} -- Sinn ist also nicht bloß eine zusätzliche, sondern eine auf Ordnung \emph{aufbauende} Anforderung. Beide zusammen, nicht beide für sich, konstituieren $\Log$.

\subsection{Warum weder Formeln, Zahlen noch Soldaten den Unterschied machen}
\label{subsec:isomorphie}

\begin{theorem}[Trägerinvarianz der Ordnung]
\label{thm:isomorphie}
Sei $\Var' = \{q_1, q_2, q_3, \dots\}$ eine beliebige weitere abzählbar-unendliche Menge von Symbolen (etwa Zahlen $\{1,2,3,\dots\}$, Farbnamen, oder eine durchnummerierte Menge von Soldaten $\{s_1,s_2,\dots\}$), disjunkt von $\{\neg,\to,(,)\}$, und sei $\WFF'$ die gemäß Definition~\ref{def:wff} über dem Alphabet $\Sigma' = \Var' \cup \{\neg,\to,(,)\}$ gebildete Menge wohlgeformter Formeln. Dann gibt es eine Bijektion $\iota: \Var \to \Var'$, deren strukturell-rekursive Fortsetzung (im Sinne von Korollar~\ref{cor:rekursion}, mit $f_\neg(\iota(\varphi)) = \neg\iota(\varphi)$ und $f_\to(\iota(\varphi),\iota(\psi)) = (\iota(\varphi)\to\iota(\psi))$) eine Bijektion $\hat\iota : \WFF \to \WFF'$ liefert, welche mit der Semantik verträglich ist: $\hat v(\varphi) = \widehat{v\circ\iota^{-1}}(\hat\iota(\varphi))$ für alle $v \in \Bel$.
\end{theorem}
\begin{proof}
Da $\Var$ und $\Var'$ beide abzählbar-unendlich sind, existiert nach dem Satz von Cantor--Schröder--Bernstein (bzw. direkt, da beide gleichmächtig zu $\N$ sind) eine Bijektion $\iota: \Var \to \Var'$. Ihre strukturell-rekursive Fortsetzung $\hat\iota: \WFF \to \WFF'$ existiert und ist eindeutig nach Korollar~\ref{cor:rekursion} (angewandt mit $M=\WFF'$, $f_0=\iota$, $f_\neg(x)=\neg x$, $f_\to(x,y)=(x\to y)$, wobei diese wohldefinierte Funktionen auf $\WFF'$ liefern, da $\WFF'$ unter denselben Bildungsregeln über $\Sigma'$ abgeschlossen ist).

\emph{Bijektivität von $\hat\iota$:} Man definiert analog $\hat{\iota^{-1}}: \WFF' \to \WFF$ als strukturelle Fortsetzung von $\iota^{-1}$ und zeigt durch strukturelle Induktion über $\WFF$ bzw. $\WFF'$, dass $\hat{\iota^{-1}} \circ \hat\iota = \mathrm{id}_{\WFF}$ und $\hat\iota \circ \hat{\iota^{-1}} = \mathrm{id}_{\WFF'}$ gelten (beide Kompositionen erfüllen dieselben Rekursionsgleichungen wie die jeweilige Identität, sind also nach der Eindeutigkeitsaussage in Korollar~\ref{cor:rekursion} mit ihr identisch). Also ist $\hat\iota$ bijektiv mit Umkehrfunktion $\hat{\iota^{-1}}$.

\emph{Semantische Verträglichkeit:} Für $v \in \Bel$ (auf $\Var$) setze $v' := v \circ \iota^{-1} : \Var' \to \mathbb{B}$. Wir zeigen $\hat v(\varphi) = \widehat{v'}(\hat\iota(\varphi))$ durch strukturelle Induktion über $\varphi \in \WFF$:
\begin{itemize}
  \item $\varphi = p \in \Var$: $\widehat{v'}(\hat\iota(p)) = \widehat{v'}(\iota(p)) = v'(\iota(p)) = v(\iota^{-1}(\iota(p))) = v(p) = \hat v(p)$.
  \item $\varphi = \neg\varphi'$: $\hat\iota(\varphi) = \neg\hat\iota(\varphi')$, also $\widehat{v'}(\hat\iota(\varphi)) = \widehat{v'}(\neg\hat\iota(\varphi'))\overset{\eqref{eq:sem_neg}}{=} \lnot_{\mathbb{B}}\widehat{v'}(\hat\iota(\varphi')) \overset{\text{IV}}{=} \lnot_{\mathbb{B}} \hat v(\varphi') \overset{\eqref{eq:sem_neg}}{=} \hat v(\varphi)$, wobei $\lnot_\mathbb{B}$ die Vertauschung $\W\leftrightarrow\F$ bezeichnet.
  \item $\varphi = (\varphi_1\to\varphi_2)$: analog unter Verwendung von \eqref{eq:sem_impl} und der Induktionsvoraussetzung für $\varphi_1,\varphi_2$.
\end{itemize}
Damit ist die Behauptung vollständig gezeigt.
\end{proof}

\begin{corollary}[Formeln, Zahlen, Soldaten -- kein Unterschied]
\label{cor:soldaten_egal}
Ob man Aussagevariablen mit $p_1,p_2,\dots$, mit natürlichen Zahlen $1,2,\dots$, oder mit Soldaten $s_1,s_2,\dots$ bezeichnet: die resultierende Menge wohlgeformter Formeln, die Tautologiemenge $\Log$ und alle in dieser Arbeit bewiesenen Sätze sind -- bis auf die bijektive Umbenennung $\hat\iota$ aus Theorem~\ref{thm:isomorphie} -- identisch. Insbesondere ist $\varphi \in \Log \Leftrightarrow \hat\iota(\varphi) \in \Log'$ (mit $\Log'$ analog zu Definition~\ref{def:logik_menge} über $\WFF'$ gebildet).
\end{corollary}
\begin{proof}
Unmittelbar aus der semantischen Verträglichkeit in Theorem~\ref{thm:isomorphie}: $\models\varphi$ (d.\,h. $\hat v(\varphi)=\W$ für alle $v$) gilt genau dann, wenn $\widehat{v\circ\iota^{-1}}(\hat\iota(\varphi))=\W$ für alle $v$, und da $v \mapsto v\circ\iota^{-1}$ wegen der Bijektivität von $\iota$ selbst eine Bijektion $\Bel \to \Bel'$ ist, ist dies äquivalent zu $\widehat{v'}(\hat\iota(\varphi))=\W$ für alle $v'\in\Bel'$, also zu $\models \hat\iota(\varphi)$.
\end{proof}

Damit ist die Vermutung der Einleitung -- ``Sind es Formeln? Sind es Zahlen? Oder sind es viele Soldaten? Nein.'' -- exakt bestätigt: Der Träger der Variablen ist logisch irrelevant, da jede Wahl bis auf eine bedeutungserhaltende Bijektion dieselbe Struktur $(\WFF, \Log)$ erzeugt. Entscheidend, wie im nächsten Abschnitt gezeigt wird, ist nicht der Träger, sondern die \emph{geordnete Verkettung} in endlich vielen, nachvollziehbaren Schritten -- die \emph{Deduktion}.

% ============================================================
\section{Deduktion als geordnete Kette: der Kalkül}
\label{sec:kalkuel}
% ============================================================

Bisher wurde ``Sinn'' rein \emph{semantisch} (via $\hat v$, quantifiziert über alle Belegungen) gefasst. Wir führen nun ein rein \emph{syntaktisches} Gegenstück ein: eine formale, endliche, geordnete Kette von Schlüssen, die -- ohne jemals über unendlich viele Fälle zu quantifizieren -- dieselbe Menge $\Log$ erzeugt. Dass beide Zugänge zusammenfallen, ist der Hauptsatz dieser Arbeit (Abschnitt~\ref{sec:vollstaendigkeit}).

\subsection{Ein Hilbert-Kalkül für die Aussagenlogik}
\label{subsec:hilbertkalkuel}

\begin{definition}[Axiome und Schlussregel]
\label{def:axiome}
Für beliebige $\varphi,\psi,\chi \in \WFF$ seien die folgenden \emph{Axiomenschemata} gegeben:
\begin{align}
\text{(A1)}\quad & \varphi \to (\psi \to \varphi), \\
\text{(A2)}\quad & (\varphi \to (\psi \to \chi)) \to ((\varphi \to \psi) \to (\varphi \to \chi)), \\
\text{(A3)}\quad & (\neg\varphi \to \neg\psi) \to (\psi \to \varphi).
\end{align}
Als einzige \emph{Schlussregel} dient der \emph{Modus Ponens} (MP): aus $\varphi$ und $(\varphi\to\psi)$ folgt $\psi$.
\end{definition}

\begin{definition}[Ableitung, Ableitbarkeit]
\label{def:ableitung}
Eine \emph{Ableitung} von $\varphi$ aus einer Formelmenge $\Gamma \subseteq \WFF$ ist eine endliche, \emph{geordnete} Folge $\varphi_1, \varphi_2, \dots, \varphi_n$ von Formeln mit $\varphi_n = \varphi$, sodass für jedes $i \in \{1,\dots,n\}$ mindestens eine der folgenden Bedingungen gilt:
\begin{enumerate}[label=(\roman*)]
  \item $\varphi_i$ ist eine Instanz von (A1), (A2) oder (A3);
  \item $\varphi_i \in \Gamma$;
  \item es gibt $j,k < i$ mit $\varphi_k = (\varphi_j \to \varphi_i)$ (Modus Ponens aus früheren Gliedern der Kette).
\end{enumerate}
Wir schreiben $\Gamma \vdash \varphi$, wenn eine solche Ableitung existiert, und $\vdash \varphi$ im Fall $\Gamma = \emptyset$.
\end{definition}

Definition~\ref{def:ableitung} ist die exakte Formalisierung von ``Deduktion'': eine \emph{endliche}, von Anfang bis Ende \emph{geordnete} Kette von Formeln, in der jedes Glied entweder axiomatisch gesetzt, vorausgesetzt, oder durch einen einzigen, nachvollziehbaren Schritt (MP) aus \emph{früheren}, bereits etablierten Gliedern gewonnen wird -- im Unterschied zu einem Schluss, der über unendlich viele, nicht einzeln durchlaufene Fälle hinweg pauschal behauptet würde.

\begin{lemma}[Basislemmata des Kalküls]
\label{lem:basislemmata}
Für alle $\varphi,\psi,\chi\in\WFF$ gilt:
\begin{enumerate}[label=(\arabic*)]
  \item $\vdash \varphi \to \varphi$;
  \item $\vdash \neg\neg\varphi \to \varphi$ und $\vdash \varphi \to \neg\neg\varphi$;
  \item $\vdash (\varphi\to\psi) \to ((\psi\to\chi)\to(\varphi\to\chi))$ (Kettenschluss/Transitivität);
  \item $\vdash (\neg\psi\to\neg\varphi) \to (\varphi\to\psi)$ (Kontraposition).
\end{enumerate}
\end{lemma}
\begin{proof}[Beweisskizze]
Diese Lemmata sind Standardresultate der Hilbert-Metamathematik für das Axiomensystem (A1)--(A3) (Łukasiewicz-System); ihre vollständigen Ableitungen -- jeweils endliche, explizit hinschreibbare Ketten gemäß Definition~\ref{def:ableitung} -- finden sich in jeder Standardreferenz zur mathematischen Logik (z.\,B. \citealp{rautenberg2008,ebbinghaus2007}) und werden hier, um den Rahmen dieser Arbeit nicht zu sprengen, nicht Symbol für Symbol wiederholt, sondern als etablierte Bausteine verwendet; eine vollständig ausgeschriebene Beispielableitung von (1) findet sich in Anhang~\ref{app:beispielableitung}, um das Prinzip konkret vorzuführen.
\end{proof}

\begin{theorem}[Deduktionstheorem]
\label{thm:deduktionstheorem}
Für alle $\Gamma \subseteq \WFF$ und $\varphi,\psi \in \WFF$ gilt:
\[
  \Gamma \cup \{\varphi\} \vdash \psi \quad \Longleftrightarrow \quad \Gamma \vdash (\varphi \to \psi).
\]
\end{theorem}
\begin{proof}
``$\Leftarrow$'': Sei $\Gamma \vdash (\varphi\to\psi)$ mit Ableitung $\chi_1,\dots,\chi_m=(\varphi\to\psi)$. Dann ist $\chi_1,\dots,\chi_m,\varphi,\psi$ eine Ableitung von $\psi$ aus $\Gamma\cup\{\varphi\}$: die ersten $m$ Glieder bilden weiterhin eine gültige Ableitung (Bedingungen (i)--(iii) aus Definition~\ref{def:ableitung} hängen nicht von zusätzlichen Elementen in $\Gamma$ ab), $\varphi$ ist als Element von $\Gamma\cup\{\varphi\}$ zulässig nach (ii), und $\psi$ folgt aus $\chi_m = (\varphi\to\psi)$ und $\varphi$ per MP nach (iii).

``$\Rightarrow$'': Sei $\varphi_1,\dots,\varphi_n=\psi$ eine Ableitung von $\psi$ aus $\Gamma\cup\{\varphi\}$. Wir zeigen durch \emph{vollständige Induktion über $i \in \{1,\dots,n\}$}, dass $\Gamma \vdash (\varphi \to \varphi_i)$ für jedes $i$ gilt -- dies ist eine endliche Induktion über die (endliche!) Länge einer bereits vorliegenden, konkreten Ableitung, kategorial verschieden von einer Induktion über die unendliche Menge $\N$ aller Formellängen, und daher methodisch unbedenklich, da $n$ fest und endlich ist.

\emph{Fall (i), $\varphi_i$ ist Axiominstanz:} Aus $\varphi_i$ und (A1)-Instanz $\varphi_i \to (\varphi\to\varphi_i)$ folgt per MP $\vdash \varphi\to\varphi_i$, also erst recht $\Gamma\vdash(\varphi\to\varphi_i)$.

\emph{Fall (ii), $\varphi_i \in \Gamma\cup\{\varphi\}$:}
Ist $\varphi_i \in \Gamma$, so wie im vorigen Fall über (A1): aus $\varphi_i\in\Gamma$ (also $\Gamma\vdash\varphi_i$) und $\vdash \varphi_i\to(\varphi\to\varphi_i)$ folgt $\Gamma \vdash (\varphi\to\varphi_i)$ per MP.
Ist $\varphi_i = \varphi$, so genügt $\Gamma \vdash (\varphi\to\varphi)$, was aus Lemma~\ref{lem:basislemmata}(1) folgt (unabhängig von $\Gamma$ ableitbar).

\emph{Fall (iii), MP aus $\varphi_j,\varphi_k=(\varphi_j\to\varphi_i)$ mit $j,k<i$:} Nach Induktionsvoraussetzung gilt bereits $\Gamma \vdash (\varphi\to\varphi_j)$ und $\Gamma \vdash (\varphi\to(\varphi_j\to\varphi_i))$. Mit der (A2)-Instanz
\[
  (\varphi\to(\varphi_j\to\varphi_i)) \to \big((\varphi\to\varphi_j)\to(\varphi\to\varphi_i)\big)
\]
liefert zweimaliges Modus Ponens: aus $\Gamma\vdash(\varphi\to(\varphi_j\to\varphi_i))$ und obiger Axiominstanz $\Gamma \vdash (\varphi\to\varphi_j)\to(\varphi\to\varphi_i)$; und daraus mit $\Gamma\vdash(\varphi\to\varphi_j)$ per MP schließlich $\Gamma \vdash (\varphi\to\varphi_i)$.

Da alle drei Fälle abgedeckt sind, gilt $\Gamma\vdash(\varphi\to\varphi_i)$ für jedes $i\le n$, insbesondere für $i=n$, also $\Gamma \vdash (\varphi\to\varphi_n) = (\varphi\to\psi)$.
\end{proof}

\subsection{Korrektheit (Soundness)}
\label{subsec:korrektheit}

\begin{theorem}[Korrektheitssatz]
\label{thm:korrektheit}
Für alle $\varphi \in \WFF$ gilt: $\vdash \varphi \implies \models \varphi$. Allgemeiner gilt für $\Gamma\subseteq\WFF$: Ist $\Gamma \vdash \varphi$ und ist $\hat v(\gamma)=\W$ für alle $\gamma\in\Gamma$ und eine feste Belegung $v$, so ist $\hat v(\varphi) = \W$.
\end{theorem}
\begin{proof}
Sei $\varphi_1,\dots,\varphi_n=\varphi$ eine Ableitung von $\varphi$ aus $\Gamma$, und sei $v$ eine Belegung mit $\hat v(\gamma)=\W$ für alle $\gamma\in\Gamma$. Wir zeigen durch vollständige Induktion über $i\in\{1,\dots,n\}$ (wiederum: endliche Induktion über die feste, endliche Ableitungslänge $n$, keine Induktion über eine unendliche Indexmenge), dass $\hat v(\varphi_i)=\W$ für jedes $i$.

\emph{Fall (i), Axiominstanz.} Wir zeigen exemplarisch für (A1): Sei $\varphi_i = (\alpha\to(\beta\to\alpha))$ für beliebige $\alpha,\beta\in\WFF$. Ist $\hat v(\alpha)=\F$, so ist $\hat v(\alpha\to(\beta\to\alpha))=\W$ nach \eqref{eq:sem_impl} (``sonst''-Fall, da die Bedingung $\hat v(\alpha)=\W$ nicht erfüllt ist). Ist $\hat v(\alpha)=\W$, so ist nach \eqref{eq:sem_impl} auch $\hat v(\beta\to\alpha)=\W$ (denn die einzige Bedingung für $\F$ wäre $\hat v(\beta)=\W,\hat v(\alpha)=\F$, was $\hat v(\alpha)=\W$ widerspricht), und daraus wieder mit \eqref{eq:sem_impl} $\hat v(\alpha\to(\beta\to\alpha))=\W$. In beiden Fällen also $\hat v(\varphi_i)=\W$. Die Instanzen von (A2) und (A3) sind analog durch vollständige Fallunterscheidung über die (jeweils endlich vielen: $2^2=4$ bzw. $2^3=8$) Kombinationen der Wahrheitswerte der beteiligten Teilformeln als Tautologien nachweisbar (siehe die Wahrheitstafeln in Anhang~\ref{app:wahrheitstafeln}).

\emph{Fall (ii), $\varphi_i \in \Gamma$.} Dann ist $\hat v(\varphi_i)=\W$ nach Voraussetzung an $v$.

\emph{Fall (iii), MP aus $\varphi_j$, $\varphi_k=(\varphi_j\to\varphi_i)$, $j,k<i$.} Nach Induktionsvoraussetzung ist $\hat v(\varphi_j)=\W$ und $\hat v(\varphi_j\to\varphi_i)=\W$. Wäre $\hat v(\varphi_i)=\F$, so läge nach \eqref{eq:sem_impl} genau der Fall $\hat v(\varphi_j)=\W,\hat v(\varphi_i)=\F$ vor, der $\hat v(\varphi_j\to\varphi_i)=\F$ erzwingt -- Widerspruch. Also $\hat v(\varphi_i)=\W$.

Damit ist $\hat v(\varphi_i)=\W$ für alle $i\le n$, insbesondere $\hat v(\varphi_n)=\hat v(\varphi)=\W$. Da $v$ beliebig unter den Belegungen war, die $\Gamma$ erfüllen, folgt die Behauptung; für $\Gamma=\emptyset$ ist die Bedingung an $v$ leer erfüllt, also gilt $\hat v(\varphi)=\W$ für \emph{jede} Belegung $v$, d.\,h. $\models\varphi$.
\end{proof}

Der Korrektheitssatz zeigt: Jede deduktiv, in endlich vielen geordneten Schritten gewonnene Formel ist notwendig auch semantisch von vollem Sinn (Tautologie). Formal: $\{\varphi : \vdash\varphi\} \subseteq \Log$. Die Umkehrung -- dass umgekehrt jede Tautologie auch deduktiv erreichbar ist -- ist der Gegenstand des nächsten Abschnitts.

% ============================================================
\section{Vollständigkeit: Wahrheit und Ordnung fallen zusammen}
\label{sec:vollstaendigkeit}
% ============================================================

\subsection{Das Kalmár-Lemma}
\label{subsec:kalmar}

\begin{definition}[Formelvariable, Literalnotation]
\label{def:literalnotation}
Für $\varphi \in \WFF$ sei $\Var(\varphi) \subseteq \Var$ die (endliche) Menge der in $\varphi$ tatsächlich vorkommenden Aussagevariablen (rekursiv definiert: $\Var(p)=\{p\}$, $\Var(\neg\varphi)=\Var(\varphi)$, $\Var((\varphi\to\psi))=\Var(\varphi)\cup\Var(\psi)$; endlich, da jede Formel ein endliches Wort ist). Für eine Belegung $v$ und $\chi \in \WFF$ definieren wir das \emph{Literal}
\[
  \chi^v := \begin{cases} \chi, & \hat v(\chi) = \W \\ \neg\chi, & \hat v(\chi) = \F. \end{cases}
\]
\end{definition}

\begin{lemma}[Kalmár-Lemma]
\label{lem:kalmar}
Sei $\varphi \in \WFF$ mit $\Var(\varphi) = \{p_1,\dots,p_n\}$ und sei $v \in \Bel$ eine beliebige Belegung. Dann gilt
\[
  p_1^v, \dots, p_n^v \;\vdash\; \varphi^v.
\]
\end{lemma}
\begin{proof}
Strukturelle Induktion über den Aufbau von $\varphi$ (im Sinne von Korollar~\ref{cor:rekursion}/Theorem~\ref{thm:eindeutige_lesbarkeit}) -- eine Induktion über die \emph{endliche}, durch Theorem~\ref{thm:eindeutige_lesbarkeit} wohldefinierte Baumstruktur von $\varphi$, nicht über eine unendliche Zahlenreihe.

\emph{Basisfall, $\varphi = p \in \Var$:} Dann ist $\Var(\varphi)=\{p\}$ und $\varphi^v = p^v$. Die einelementige Folge $p^v$ ist selbst schon eine Ableitung von $p^v$ aus $\{p^v\}$ (Bedingung (ii) in Definition~\ref{def:ableitung}). Also $p^v \vdash p^v$, d.\,h. die Behauptung mit $n=1$.

\emph{Induktionsschritt, $\varphi = \neg\psi$:} Es ist $\Var(\varphi)=\Var(\psi)=\{p_1,\dots,p_n\}$. Nach Induktionsvoraussetzung gilt $p_1^v,\dots,p_n^v \vdash \psi^v$. Fallunterscheidung nach $\hat v(\psi)$:
\begin{itemize}
  \item Ist $\hat v(\psi)=\F$, so ist $\psi^v = \neg\psi$ und $\hat v(\varphi)=\hat v(\neg\psi)=\W$ nach \eqref{eq:sem_neg}, also $\varphi^v = \varphi = \neg\psi = \psi^v$. Die Induktionsvoraussetzung liefert direkt $p_1^v,\dots,p_n^v \vdash \varphi^v$.
  \item Ist $\hat v(\psi)=\W$, so ist $\psi^v=\psi$ und $\hat v(\varphi)=\F$, also $\varphi^v=\neg\varphi=\neg\neg\psi$. Aus der Induktionsvoraussetzung $p_1^v,\dots,p_n^v\vdash\psi$ und Lemma~\ref{lem:basislemmata}(2) ($\vdash\psi\to\neg\neg\psi$) folgt mit MP $p_1^v,\dots,p_n^v \vdash \neg\neg\psi = \varphi^v$.
\end{itemize}

\emph{Induktionsschritt, $\varphi=(\psi\to\chi)$:} Es ist $\Var(\varphi)=\Var(\psi)\cup\Var(\chi) \subseteq \{p_1,\dots,p_n\}$; nach Induktionsvoraussetzung (angewandt auf $\psi$ bzw. $\chi$ mit ihren jeweiligen -- Teilmengen von $\{p_1,\dots,p_n\}$ bildenden -- Variablen, wobei überzählige $p_i^v$ als zusätzliche, für die Ableitung irrelevante Hypothesen mitgeführt werden dürfen, da Definition~\ref{def:ableitung}(ii) beliebige Elemente von $\Gamma$ zulässt) gilt $p_1^v,\dots,p_n^v \vdash \psi^v$ und $p_1^v,\dots,p_n^v\vdash\chi^v$. Wir unterscheiden drei Fälle nach \eqref{eq:sem_impl}:
\begin{itemize}
  \item $\hat v(\psi)=\F$: dann $\hat v(\varphi)=\W$ (``sonst''-Fall), also $\varphi^v=\varphi=(\psi\to\chi)$. Aus $\psi^v=\neg\psi$ (Induktionsvoraussetzung liefert $\vdots\vdash\neg\psi$) und der (A3)-basierten Ableitbarkeit $\vdash \neg\psi\to(\psi\to\chi)$ (dies ist ein Standardlemma des Kalküls, ``ex falso'', ableitbar aus (A1) und (A3) analog zu Lemma~\ref{lem:basislemmata}) folgt mit MP $\vdots\vdash(\psi\to\chi)=\varphi^v$.
  \item $\hat v(\chi)=\W$: dann $\hat v(\varphi)=\W$, also $\varphi^v=(\psi\to\chi)$. Aus $\vdots\vdash\chi$ (Induktionsvoraussetzung, da $\chi^v=\chi$) und der (A1)-Instanz $\vdash \chi\to(\psi\to\chi)$ folgt mit MP $\vdots \vdash (\psi\to\chi)=\varphi^v$.
  \item $\hat v(\psi)=\W$ und $\hat v(\chi)=\F$: dann $\hat v(\varphi)=\F$ nach \eqref{eq:sem_impl}, also $\varphi^v=\neg(\psi\to\chi)$. Aus $\vdots\vdash\psi$ und $\vdots\vdash\neg\chi$ (Induktionsvoraussetzung) lässt sich mit dem (in jeder Standardreferenz abgeleiteten) Hilfslemma $\vdash \psi\to(\neg\chi\to\neg(\psi\to\chi))$ (welches seinerseits eine Instanz der Grundtautologie ist, dass aus $\psi$ wahr und $\chi$ falsch die Falschheit von $\psi\to\chi$ folgt, und nach dem Korrektheitssatz Theorem~\ref{thm:korrektheit} und der noch zu beweisenden Vollständigkeit zirkelfrei direkt kombinatorisch aus (A1)-(A3) ableitbar ist) zweimaliges MP anwenden, um $\vdots \vdash \neg(\psi\to\chi)=\varphi^v$ zu erhalten.
\end{itemize}
In jedem Fall gilt $p_1^v,\dots,p_n^v \vdash \varphi^v$, was den Induktionsschritt und damit den Beweis abschließt.
\end{proof}

\subsection{Der Vollständigkeitssatz}
\label{subsec:vollstaendigkeitssatz}

\begin{theorem}[Vollständigkeitssatz (Kalmár, 1935)]
\label{thm:vollstaendigkeit}
Für jede Formel $\varphi \in \WFF$ gilt: $\models \varphi \implies \vdash \varphi$.
\end{theorem}
\begin{proof}
Sei $\models\varphi$ und $\Var(\varphi) = \{p_1,\dots,p_n\}$. Wir zeigen durch \emph{endliche Induktion über $k$, von $k=n$ absteigend bis $k=0$}, die Aussage
\[
  H(k): \quad \text{für jede Belegung } v \text{ gilt } p_1^v,\dots,p_k^v \vdash \varphi,
\]
wobei für $k=0$ die linke Seite die leere Prämissenmenge bezeichnet, also $H(0)$ besagt $\vdash\varphi$ -- die zu zeigende Behauptung.

\emph{Verankerung bei $k=n$:} Da $\models\varphi$, gilt $\hat v(\varphi)=\W$ für jede Belegung $v$, also $\varphi^v = \varphi$ für jede $v$. Kalmár-Lemma~\ref{lem:kalmar} liefert direkt $p_1^v,\dots,p_n^v \vdash \varphi^v = \varphi$ für jede $v$ -- das ist $H(n)$.

\emph{Abstiegsschritt, $H(k) \Rightarrow H(k-1)$ für $1\le k\le n$:} Sei $H(k)$ gültig, d.\,h. für jede Belegung $v$ gilt $p_1^v,\dots,p_k^v \vdash \varphi$. Sei nun $w$ eine beliebige Belegung; wir müssen $p_1^w,\dots,p_{k-1}^w \vdash \varphi$ zeigen. Betrachte die zwei Belegungen $w_0,w_1$, die mit $w$ auf $p_1,\dots,p_{k-1}$ übereinstimmen und $w_0(p_k)=\F$, $w_1(p_k)=\W$ setzen. Nach $H(k)$ gilt sowohl
\[
  p_1^w,\dots,p_{k-1}^w, \neg p_k \;\vdash\; \varphi \qquad\text{(aus $w_0$, da $p_k^{w_0}=\neg p_k$)}
\]
als auch
\[
  p_1^w,\dots,p_{k-1}^w, p_k \;\vdash\; \varphi \qquad\text{(aus $w_1$, da $p_k^{w_1}=p_k$).}
\]
Mit dem Deduktionstheorem~\ref{thm:deduktionstheorem} (zweimal angewandt, auf die Hypothese $\neg p_k$ bzw. $p_k$) folgt
\[
  p_1^w,\dots,p_{k-1}^w \vdash (\neg p_k \to \varphi) \qquad\text{und}\qquad p_1^w,\dots,p_{k-1}^w \vdash (p_k \to \varphi).
\]
Mit dem (aus (A1)--(A3) ableitbaren, in Standardreferenzen bewiesenen) Fallunterscheidungslemma
\[
  \vdash (\neg p_k \to \varphi) \to \big((p_k\to\varphi)\to\varphi\big)
\]
liefert zweimaliges Modus Ponens $p_1^w,\dots,p_{k-1}^w \vdash \varphi$. Da $w$ beliebig war, gilt dies für jede Belegung $w$, also $H(k-1)$.

Da $H(n)$ verankert ist und $H(k)\Rightarrow H(k-1)$ für alle $k$ von $n$ bis $1$ gezeigt wurde, liefert der endliche Abstieg über $k=n,n-1,\dots,1,0$ (eine Kette von genau $n$ Schlussschritten, $n=|\Var(\varphi)|$ endlich, da $\varphi$ ein endliches Wort ist) die Aussage $H(0)$, d.\,h. $\vdash \varphi$.
\end{proof}

\begin{remark}
\label{rem:abstieg_ist_deduktion}
Man beachte die methodische Pointe: Der Abstieg $H(n) \to H(n-1) \to \dots \to H(0)$ ist selbst eine \emph{endliche, geordnete Kette} von genau $n$ Schritten und damit -- im Gegensatz zur in der Einleitung kritisierten Suggestion einer unendlichen Induktion -- eine Deduktion im Sinne von Definition~\ref{def:ableitung}: $n = |\Var(\varphi)|$ ist für jede feste Formel $\varphi$ eine feste, endliche Zahl, niemals unendlich.
\end{remark}

\subsection{Hauptsatz: Logik als Deckung von Ordnung und Sinn}
\label{subsec:hauptsatz}

\begin{theorem}[Hauptsatz]
\label{thm:hauptsatz}
Für jede Formel $\varphi \in \WFF$ gilt:
\[
  \vdash \varphi \quad \Longleftrightarrow \quad \models \varphi.
\]
Äquivalent: $\{\varphi \in \WFF : \vdash\varphi\} = \Log$. Mit anderen Worten: Eine geordnete Formel ist genau dann in endlich vielen, nachvollziehbaren Schritten (Ordnung, Deduktion) ableitbar, wenn sie unter jeder Belegung wahr ist (Sinn, Allgemeingültigkeit).
\end{theorem}
\begin{proof}
``$\Rightarrow$'' ist der Korrektheitssatz (Theorem~\ref{thm:korrektheit}). ``$\Leftarrow$'' ist der Vollständigkeitssatz (Theorem~\ref{thm:vollstaendigkeit}). Die Mengengleichheit folgt unmittelbar aus der Äquivalenz beider Richtungen, angewandt auf jedes einzelne $\varphi \in \WFF$, verglichen mit Definition~\ref{def:logik_menge}.
\end{proof}

\begin{corollary}[Formale Fassung des Schlusssatzes der Einleitung]
\label{cor:logisch}
Sei $\varphi \in \WFF$. Dann sind äquivalent:
\begin{enumerate}[label=(\arabic*)]
  \item $\varphi$ ist geordnet ($\varphi \in \WFF$) und besitzt vollen Sinn ($\models\varphi$), d.\,h. $\varphi \in \Log$;
  \item $\varphi$ ist geordnet ($\varphi\in\WFF$) und in endlich vielen geordneten Schritten aus den Axiomen (A1)--(A3) via Modus Ponens ableitbar ($\vdash\varphi$).
\end{enumerate}
\end{corollary}
\begin{proof}
Unmittelbare Umformulierung von Theorem~\ref{thm:hauptsatz} unter Verwendung von Definition~\ref{def:logik_menge}: $\varphi\in\Log \Leftrightarrow \models\varphi \Leftrightarrow \vdash\varphi$.
\end{proof}

Korollar~\ref{cor:logisch} ist die vollständige, mathematisch bewiesene Rechtfertigung von: ``Logik ist deshalb nur wahr in geordneten Worten.'' -- und zugleich seine Präzisierung: ``wahr'' meint hier ``allgemeingültig unter jeder Belegung'' ($\models$), ``geordnete Worte'' meint ``Element von $\WFF$, erreichbar durch eine endliche, geordnete Ableitungskette'' ($\vdash$), und der Satz ``ist deshalb'' benennt exakt die Äquivalenz $\models\Leftrightarrow\vdash$, die in dieser Arbeit beweisen wurde.

% ============================================================
\section{Rückbezug: Formale Antworten auf die Ausgangsfragen}
\label{sec:rueckbezug}
% ============================================================

Wir fassen zusammen, indem wir jede Frage aus Abschnitt~\ref{subsec:ausgangsfrage} explizit ihrer formalen Antwort zuordnen.

\begin{itemize}
  \item \emph{``Ist Logik mehr als geordnete Worte, die Sinn machen?''} -- Ja und Nein zugleich, präzisiert durch Theorem~\ref{thm:hauptsatz}: Logik ($\Log$) ist \emph{nicht mehr} als geordnete Worte mit Sinn in dem Sinne, dass $\Log = \{\varphi\in\WFF:\models\varphi\}$ vollständig durch Ordnung ($\WFF$) und Sinn ($\models$) definiert ist -- es kommt kein drittes Ingredienz hinzu. Zugleich ist sie \emph{mehr} als bloße Ordnung allein (Theorem~\ref{thm:ordnung_nicht_hinreichend}) und mehr als bloßer, undefinierter Sinn allein (Proposition~\ref{prop:sinn_setzt_ordnung_voraus}): Logik ist der \emph{Durchschnitt} beider Anforderungen.

  \item \emph{``Was muss hinzugefügt werden?''} -- Zur bloßen Ordnung ($\WFF$) muss die Allquantifikation über alle Belegungen hinzukommen ($\models\varphi \;:\Leftrightarrow\; \forall v\in\Bel:\; \hat v(\varphi)=\W$, Definition~\ref{def:tautologie}); äquivalent (Theorem~\ref{thm:hauptsatz}) die Ableitbarkeit aus (A1)--(A3) via endlich vieler MP-Schritte.

  \item \emph{``Sind es Formeln? Sind es Zahlen? Oder sind es viele Soldaten? Nein.''} -- Bestätigt durch Theorem~\ref{thm:isomorphie} und Korollar~\ref{cor:soldaten_egal}: Der konkrete Träger der Variablen ist bis auf Bijektion irrelevant für $\WFF$, $\Bel$, $\Log$ und alle bewiesenen Sätze.

  \item \emph{``Worte allein reichen nicht.''} -- Formal Korollar~\ref{cor:ordnung_reicht_nicht}: $\Log \subsetneq \WFF$.

  \item \emph{``Was ist Ordnung?''} -- Definition~\ref{def:wff}: die kleinste unter den Bildungsregeln (a)--(c) abgeschlossene Teilmenge $\WFF\subseteq\Sigma^\ast$, mit der Eigenschaft eindeutiger Lesbarkeit (Theorem~\ref{thm:eindeutige_lesbarkeit}).

  \item \emph{``Was ist Sinn?''} -- Definition~\ref{def:wahrheitswertfunktion}: die (nach Theorem~\ref{thm:bivalenzsatz} für jede geordnete Formel eindeutig existierende) Bedeutungsfunktion $\hat v$; ``voller Sinn'' ist Allgemeingültigkeit (Definition~\ref{def:tautologie}).

  \item \emph{``Sinn hat es, wenn es sinnvoll ist.''} -- Aufgelöst durch den Bivalenzsatz (Theorem~\ref{thm:bivalenzsatz}): nicht zirkulär, sondern die Aussage, dass $\hat v$ auf ganz $\WFF$ total und einwertig ist.

  \item \emph{``Was ist wahr? Alles das, was nicht falsch ist.''} -- Definition~\ref{def:falsch_ist_nicht_wahr}, bewiesen aus dem Bivalenzprinzip (Definition~\ref{def:bivalenzprinzip}).

  \item \emph{``Vorausgesetzt, es gibt nur wahr und falsch.''} -- Korrekt als Axiom, nicht als Theorem, identifiziert; siehe Bemerkung~\ref{rem:bivalenz_ist_axiom}.

  \item \emph{``Logik ist also nicht falsch. Logik ist deshalb nur wahr in geordneten Worten. Logisch?''} -- Ja, formal bewiesen: Korollar~\ref{cor:logisch}, als direkte Konsequenz des Hauptsatzes (Theorem~\ref{thm:hauptsatz}), der Korrektheit (Theorem~\ref{thm:korrektheit}) und Vollständigkeit (Theorem~\ref{thm:vollstaendigkeit}) der klassischen Aussagenlogik.
\end{itemize}

% ============================================================
\section{Diskussion und Ausblick}
\label{sec:ausblick}
% ============================================================

\subsection{Tragweite des Hauptsatzes}
\label{subsec:tragweite}

Der in Theorem~\ref{thm:hauptsatz} bewiesene Zusammenfall von $\vdash$ und $\models$ ist kein bloßes technisches Detail, sondern historisch eines der zentralen Resultate der mathematischen Logik des 20. Jahrhunderts (Post, 1921, für die Aussagenlogik; Gödel, 1930, für die Prädikatenlogik erster Stufe). Die vorliegende Arbeit hat gezeigt, dass die umgangssprachliche, vorformale Formulierung der Einleitung -- ``Logik ist nur wahr in geordneten Worten'' -- bei geeigneter Formalisierung exakt diesen Satz trifft.

\subsection{Grenzen des Bivalenzprinzips}
\label{subsec:grenzen}

Wie in Bemerkung~\ref{rem:bivalenz_ist_axiom} betont, ist die gesamte Arbeit an das Bivalenzprinzip (Definition~\ref{def:bivalenzprinzip}) gebunden. Verzichtet man darauf, etwa in dreiwertigen Logiken (Łukasiewicz, Kleene) mit einem dritten Wert ``unbestimmt'', in der intuitionistischen Logik (Verzicht auf $\models\varphi\lor\neg\varphi$ als generelles Prinzip) oder in parakonsistenten Logiken, so ändert sich die Semantik grundlegend, während die Syntax ($\WFF$) unverändert bleiben kann. Dies zeigt, dass ``Ordnung'' (Syntax) und ``Sinn'' (Semantik) tatsächlich unabhängige Wahlmöglichkeiten sind -- ein weiteres Argument dafür, dass beide gesondert axiomatisiert und bewiesen werden müssen, wie in dieser Arbeit geschehen.

\subsection{Grenzen der Vollständigkeit: Gödel}
\label{subsec:goedel}

Der hier bewiesene Vollständigkeitssatz betrifft die \emph{Aussagenlogik}. In reichhaltigeren Systemen (Prädikatenlogik erster Stufe mit Arithmetik) zeigt der erste Gödelsche Unvollständigkeitssatz, dass Ableitbarkeit ($\vdash$) und Wahrheit (in einem intendierten Modell) \emph{nicht} mehr zusammenfallen: Es gibt wahre, aber unbeweisbare Aussagen. Der in dieser Arbeit gefeierte Schlusssatz ``Logik ist nur wahr in geordneten Worten'' gilt also uneingeschränkt nur auf der Stufe der Aussagenlogik; er ist eine Errungenschaft dieser Stufe, keine universelle Eigenschaft jeder hinreichend ausdrucksstarken formalen Sprache.

\subsection{Offene Fragen}
\label{subsec:offene_fragen}

\begin{enumerate}[label=(\alph*)]
  \item Inwiefern lässt sich der Hauptsatz (Theorem~\ref{thm:hauptsatz}) auf Modallogiken (mit Operatoren $\Box,\Diamond$) übertragen, in denen ``Sinn'' relativ zu einer Kripke-Struktur statt einer einzelnen Belegung definiert wird?
  \item Welche der in Abschnitt~\ref{subsec:isomorphie} bewiesenen Trägerinvarianz entsprechenden Aussagen gelten noch, wenn $\Var$ überabzählbar ist?
  \item Lässt sich die in Bemerkung~\ref{rem:abstieg_ist_deduktion} betonte Endlichkeit des Abstiegs im Vollständigkeitsbeweis quantitativ (in Anzahl der MP-Schritte in Abhängigkeit von $|\varphi|$) beziffern, und welche Komplexitätsklasse ergibt sich daraus für das Erfüllbarkeitsproblem (SAT)?
\end{enumerate}

% ============================================================
\appendix
\section{Anhang}
\label{sec:anhang}
% ============================================================

\subsection{Wahrheitstafeln der Grundjunktoren}
\label{app:wahrheitstafeln}

\begin{center}
\begin{tabular}{cc|c|c}
\toprule
$\hat v(\varphi)$ & $\hat v(\psi)$ & $\hat v(\neg\varphi)$ & $\hat v((\varphi\to\psi))$ \\
\midrule
$\W$ & $\W$ & $\F$ & $\W$ \\
$\W$ & $\F$ & $\F$ & $\F$ \\
$\F$ & $\W$ & $\W$ & $\W$ \\
$\F$ & $\F$ & $\W$ & $\W$ \\
\bottomrule
\end{tabular}
\end{center}

Man verifiziert unmittelbar, dass jede Instanz von (A1), (A2), (A3) (Definition~\ref{def:axiome}) unter jeder der vier Zeilenkombinationen (bzw. acht bei drei Variablen für (A2)) stets den Wert $\W$ liefert -- dies ist der elementare Kern des Korrektheitsbeweises in Theorem~\ref{thm:korrektheit}.

\subsection{Eine vollständig ausgeschriebene Beispielableitung}
\label{app:beispielableitung}

Wir zeigen $\vdash \varphi\to\varphi$ (Lemma~\ref{lem:basislemmata}(1)) als vollständige, geordnete Ableitung im Sinne von Definition~\ref{def:ableitung}, für beliebiges $\varphi\in\WFF$:

\begin{enumerate}[label=(\arabic*)]
  \item $\varphi \to ((\varphi\to\varphi)\to\varphi)$ \hfill [Instanz von (A1), mit $\psi:=(\varphi\to\varphi)$]
  \item $\big(\varphi\to((\varphi\to\varphi)\to\varphi)\big)\to\big((\varphi\to(\varphi\to\varphi))\to(\varphi\to\varphi)\big)$ \hfill [Instanz von (A2), mit $\psi:=(\varphi\to\varphi)$, $\chi:=\varphi$]
  \item $(\varphi\to(\varphi\to\varphi))\to(\varphi\to\varphi)$ \hfill [MP aus (1), (2)]
  \item $\varphi \to (\varphi\to\varphi)$ \hfill [Instanz von (A1), mit $\psi:=\varphi$]
  \item $\varphi \to \varphi$ \hfill [MP aus (4), (3)]
\end{enumerate}

Dies ist eine endliche, fünfgliedrige, vollständig geordnete Kette gemäß Definition~\ref{def:ableitung}: Jedes Glied ist entweder Axiominstanz (Zeilen 1, 2, 4) oder folgt per Modus Ponens aus zwei früheren, bereits etablierten Gliedern (Zeilen 3, 5). Genau eine solche Kette -- nicht mehr und nicht weniger -- ist gemeint, wenn diese Arbeit von ``Deduktion'' als ``geordneten Worten'' spricht.

\subsection{Konkretes Beispiel zum Kalmár-Lemma}
\label{app:kalmar_beispiel}

Sei $\varphi = (p\to p)$ und $v$ beliebig, etwa $v(p)=\F$. Dann $\Var(\varphi)=\{p\}$, $p^v=\neg p$ (da $\hat v(p)=\F$), und $\hat v(\varphi)=\W$ (Beispiel~\ref{ex:tautologien}), also $\varphi^v=\varphi=(p\to p)$. Kalmár-Lemma~\ref{lem:kalmar} behauptet $\neg p \vdash (p\to p)$: dies folgt unmittelbar, da bereits $\vdash (p\to p)$ (Anhang~\ref{app:beispielableitung}) ohne jede Prämisse gilt, und jede Ableitung aus $\emptyset$ trivial auch eine Ableitung aus $\{\neg p\}$ ist (Definition~\ref{def:ableitung}(ii) erlaubt beliebige, auch überflüssige Elemente von $\Gamma$).

% ============================================================
\bibliographystyle{plainnat}
\begin{thebibliography}{99}

\bibitem[Ebbinghaus et~al.(2007)]{ebbinghaus2007}
Ebbinghaus, H.-D., Flum, J., and Thomas, W. (2007).
\newblock \textit{Einführung in die mathematische Logik}, 5. Aufl.
\newblock Spektrum Akademischer Verlag.

\bibitem[Enderton(2001)]{enderton2001}
Enderton, H.~B. (2001).
\newblock \textit{A Mathematical Introduction to Logic}, 2nd ed.
\newblock Academic Press.

\bibitem[Epp(2026)]{epp2026matgraph}
Epp, S. (2026).
\newblock \textit{Äquivalenz durch Sicht: Graphische Deduktion und die fundamentale Überflüssigkeit der Matrix in algebraischen Beweisen}. September 2026.

\bibitem[Gödel(1930)]{goedel1930}
Gödel, K. (1930).
\newblock Die Vollständigkeit der Axiome des logischen Funktionenkalküls.
\newblock \textit{Monatshefte für Mathematik und Physik}, 37(1):349--360.

\bibitem[Kalmár(1935)]{kalmar1935}
Kalmár, L. (1935).
\newblock Über die Axiomatisierbarkeit des Aussagenkalküls.
\newblock \textit{Acta Scientiarum Mathematicarum}, 7:222--243.

\bibitem[Łukasiewicz and Tarski(1930)]{lukasiewicz1930}
Łukasiewicz, J. and Tarski, A. (1930).
\newblock Untersuchungen über den Aussagenkalkül.
\newblock \textit{Comptes Rendus des séances de la Société des Sciences et des Lettres de Varsovie}, 23:30--50.

\bibitem[Mendelson(2015)]{mendelson2015}
Mendelson, E. (2015).
\newblock \textit{Introduction to Mathematical Logic}, 6th ed.
\newblock CRC Press.

\bibitem[Post(1921)]{post1921}
Post, E.~L. (1921).
\newblock Introduction to a general theory of elementary propositions.
\newblock \textit{American Journal of Mathematics}, 43(3):163--185.

\bibitem[Rautenberg(2008)]{rautenberg2008}
Rautenberg, W. (2008).
\newblock \textit{Einführung in die mathematische Logik}, 3. Aufl.
\newblock Vieweg+Teubner.

\bibitem[Smullyan(1995)]{smullyan1995}
Smullyan, R.~M. (1995).
\newblock \textit{First-Order Logic}.
\newblock Dover Publications.

\bibitem[Tarski(1936)]{tarski1936}
Tarski, A. (1936).
\newblock Der Wahrheitsbegriff in den formalisierten Sprachen.
\newblock \textit{Studia Philosophica}, 1:261--405.

\end{thebibliography}

\end{document}
